\section{Domain-stable sequence}

We assume a SC call graph in this section.

\begin{Def}[Injective]\rm
A relation $R$ is defined to be {\em injective}
if $x R z$ and $y R z$ imply $x=y$.
\end{Def}

\begin{Def}[Composition of SC graphs]\rm
The composition $R_1\ldots R_k$ of a given SC path $(R_1,\ldots,R_k)$ is
defined as

$(f_1,(S_1,P_1)\ldots(S_k,P_k),c_{k+1},f_{k+1})$

where

$R_i=(f_i,(S_i,P_i),c_{i+1},f_{i+1})$ for all $i \in [1,k]$,

and we define

$(S_1,P_1)\ldots(S_k,P_k) = (S,P)$

where

$S = S_1\ldots S_k$ (ordinary composition of relations),

$P = (S_1 \cup P_1)\ldots (S_k \cup P_k) - S$.

We write $i(R_1\ldots R_i)j$ for $i (S \cup P) j$.

We call a composition of a SC path a {\em SC composition}.

We will often write $R$ for a SC composition.

A sequece $(R_1,\ldots,R_k)$ of SC compsitions 
is defined to be {\em composable} if
$R_i = (f_i,-,-,f_{i+1})$ for all $i \in [1,k-1]$

Note that a SC path is a special case of a composable sequence of SC compositions.

A composition of a composable sequence of SC compositions
is defined in a similar way to a composition of a SC path.

A SC composition $R$ is called {\em injective} if
the relation $S \cup P$ is injective
where
the composition $R = (-,(S,P),-,-)$.
A sequence $(R_1,\ldots,R_k)$ of SC compositions is defined to be {\em injective} if
each $R_i$ is injective for $i \in [1,k]$.
As a special case,
a SC path is defined to be {\em injective} in a similar way to 
a sequence of SC compositions.
A SC graph is defined to be {\em injective} if all edges in the graph are
injective.

For a SC composition $R$, we define

$\dom(R)$ as $\{x \ |\ R = (-,(S,P),-,-), x(S \cup P)y \}$ and

$\codom(R)$ as $\{y \ |\ R = (-,(S,P),-,-), x(S \cup P)y \}$.

For a finite set $X$, we write $|X|$ for the cardinal of $X$.

Note that $\dom(R)$ and $\codom(R)$ for a SC path $R$ is 
defined as a special case of this definition.
\end{Def}

\begin{Def}[Domain-stable (Same as def 3.3 in "Poly-exponential ...")]\rm
A sequence $(R_1,\ldots,R_k)$ of SC compositions
is defined to be {\em domain-stable}
if 
the sequence is composable, and
for all $i \in [1,k-1]$, 
we have $\dom(R_i) = \dom(R_iR_{i+1})$.
\end{Def}

\begin{Lemma}[Cardinal of domains]\label{lemma:cardinal}
If $|\dom(R)|=|\dom(RS)|$, then $(R,S)$ is domain-stable.
\end{Lemma}

{\em Proof.}
By definition of $\dom$, we have $\dom(R) \supseteq \dom(RS)$.
Since $|\dom(R)|=|\dom(RS)|$, we have $\dom(R) = \dom(RS)$.
$\Box$

\begin{Def}\rm
For a given composable sequence $R_1,\ldots,R_n$ of SC compositions,
a sequence $R'_1,\ldots,R'_l$ of SC compositions
is defined to be a {\em partial composition} of $R_1,\ldots,R_n$
with head indexes $a_1,\ldots,a_l$
if 

- there are $R_{a_1},\ldots,R_{b_1},R_{a_2},\ldots,R_{b_2},\ldots, R_{a_l},\ldots,R_{b_l}$
which is the same sequence as $R_1,\ldots,R_n$
where $a_i \le b_i$ for all $i \in [1,l]$,

- $R_i' = R_{a_i}\ldots R_{b_i}$ for all $i \in [1,l]$.
\end{Def}

\begin{Def}\rm
We define a {\em value assignment} to be a map from $N$ to $N$.
Note that we will use a value assignment $\alpha$ so that
$\alpha(i)=m$ means that the $i$-th variable has a value $m$.

We write a sequence $(v_1,\ldots,v_n)$ of natural numbers for
a value assignment $\alpha$ such that
$\alpha(i)=v_i$ for all $i \in [1,n]$ and
$\alpha(i)=0$ otherwise.

For a SC composition $R$ and value assignments $\alpha$ and $\alpha'$,
we define $\alpha R \alpha'$ to hold if for all $k,l \in N$,

- $\alpha(k) \ge \alpha'(l)$ if $k S l$, and

- $\alpha(k) > \alpha'(l)$ if $k P l$,

where
$R_i = (f,(S,P),-,g)$.
\end{Def}

Note that $\exists \alpha_2(\alpha_1 R_1 \alpha_2 R_2 \alpha_3)$ 
implies $\alpha_1 R_1R_2 \alpha_3$, but
$\alpha_1 R_1R_2 \alpha_3$ does not imply 
 $\exists \alpha_2(\alpha_1 R_1 \alpha_2 R_2 \alpha_3)$ 
since $\alpha_1 R_1 \alpha_2 R_2 \alpha_3$ implies 
$\alpha_1(1) > \alpha_2(1) > \alpha_3(1)$ 
when $1 P_1 1 P_2 1$.

\begin{Def}[Value flow]\rm
We define $\sigma$ to be a {\em value flow} along a composable sequence $\pi$ of
SC compositions 
if $\pi$ is $R_1,\ldots, R_n$,
and $\sigma$ is $\Vec v_1,\ldots,\Vec v_n$,
and 

- $\Vec v_i R_i \Vec v_{i+1}$ for all $i \in [1,n-1]$.

For a SC composition $R$ and a value assignment $\alpha$,
we define $\Val(R,\alpha) \in N$
as $\Sigma \{ \alpha(i) \ |\ i \in \dom(R) \}$.
\end{Def}

\begin{Lemma}\label{lemma:A}
if $R_1,R_2,R_3$ is composable,
and $\Vec v_1,\Vec v_2$ is along $R_1,R_2$, 
and $\Vec v_2,\Vec v_3$ is along $R_2,R_3$, 
then
$\Vec v_1,\Vec v_3$ is along $R_1R_2,R_3$.
\end{Lemma}

{\em Proof.}
Since $\exists j(i R_1 j R_2 k)$ implies $i R_1R_2 k$,
we have $\Vec v_1 (R_1R_2) \Vec v_3$.
$\Box$

\begin{Prop}[Transform to domain-stable sequence]\label{prop:domain-stable}
(1) For a composable sequence $R_1,\ldots,R_k,R$ of SC compositions,
if $R_1,\ldots,R_k$ is domain-stable
and $R_k,R$ is not domain-stable, then
the following sequcne $R_1,\ldots,R_{i-1},R'$ is
domain-stable and a partial composition of $R_1,\ldots,R_k,R$
with head redexes $1,\ldots,i-1,i$ and $R_{i},R_{i+1}\ldots R_k R$ is
not domain-stable.

- Take $i$ to be the smallest $i \in [1,k]$ among $i$'s such that
$R_i, R_{i+1} \ldots R_k R$ is not domain-stable.
It exists since the sequence $R_k,R$ is not domain-stable.
Let
\[
R' = R_i\ldots R_k R.
\]

(2) For any composable sequence $R_1,\ldots,R_k$ of SC compositions,
the following sequence $R_1',\ldots,R_l'$ 
is domain-stable and
a partial composition of 
$R_1,\ldots,R_k$ with head indexes $d_1,\ldots,d_l$.
Moreover,
if $\Vec v_1,\ldots,\Vec v_k$ is along $R_1,\ldots,R_k$,
then $\Vec v'_1,\ldots,\Vec v'_l$ is along $R'_1,\ldots,R'_l$,
where $\Vec v'_i$ is defined as $\Vec v_{d_i}$.

For each $0 < n \le k$, we define a sequence $R^{(n)}$ and 
a sequence $H^{(n)}$ of the same length by induction on $n$.
\[
R^{(1)} = (R_1), \\ 
H^{(1)} = (1).
\]
We define $R^{(n+1)}$ and $H^{(n+1)}$ from $R^{(n)}$ and $H^{(n)}$ as follows.
Let
\[
R^{(n)} = (R_1',\ldots,R_m'), \\
H^{(n)} = (n_1,\ldots,n_m).
\]
Case 1. If $R_1',\ldots,R_m',R_{n+1}$ is domain-stable,
then
\[
R^{(n+1)} = (R_1',\ldots,R_m',R_{n+1}), \\
H^{(n+1)} = (n_1,\ldots,n_m,n+1).
\]
Case 2. If $R_1',\ldots,R_m',R_{n+1}$ is not domain-stable.

By (1), there is some $i \in [1,m]$ such that
$R'_1,\ldots,R'_{i-1}, R_i'\ldots R_m'R_{n+1}$ is domain-stable
and $R_i',R_{i+1}'\ldots R_m'R_{n+1}$ is not domain-stable.
Define
\[
R^{(n+1)} = (R_1',\ldots,R_{i-1}',R_i'R'_{i+1}\ldots R'_m R_{n+1}), \\
H^{(n+1)} = (n_1,\ldots,n_{i}), \\
\]

We define $R_1',\ldots,R_l'$ as $R^{(k)}$
and define $n_1,\ldots,n_l$ as $H^{(k)}$.
\end{Prop}

{\em Proof.}
We can show $R^{(n)}$ is domain-stable for all $n \in [1,k]$
by induction on $n$.

If $\Vec v_1,\ldots,\Vec v_k$ is along $R_1,\ldots,R_k$,
then by Lemma \ref{lemma:A} $\Vec v'_1,\ldots,\Vec v'_l$ is along $R'_1,\ldots,R'_l$.
$\Box$

In this proposition,
the resulting sequence is obtained from 
maximal grouping of a given sequence of compositions 
such that a composition of elements in each group gives 
an element of a domain-stable sequence.
Here a head index is taken from the index of the head element of each group.

\begin{Lemma}\label{lemma:domain-stable-increase}
If
a composable sequence $R_1,\ldots,R_k$ of SC compositions
is domain-stable and injective, then
\[
|\dom(R_1)| \le \ldots \le |\dom(R_k)|.
\]
\end{Lemma}

{\em Proof.}
Let $R = R_i$ and $S = R_{i+1}$.

Since $R,S$ is domain-stable, we have $\dom(RS) = \dom(R)$.
Hence for each $x \in \dom(R)$ there are $y,z$ such that $x R y S z$.
Define a function $f: \dom(R) \to \dom(S)$ such that $f(x)=y$.
Since $f$ is injective, we have $|\dom(R)| \le |\dom(S)|$.
$\Box$

\begin{Lemma}[Grouping of cardinal-constant domains]\label{lemma:cardinal-constant}
For any 
domain-stable injective
sequence $R_1,\ldots,R_k$ of SC compositions,
there are $a_1,b_1,\ldots, a_l,b_l \in [1,k]$ such that
$R_{a_1},\ldots,R_{b_1},R_{a_2},\ldots,R_{b_2},\ldots,R_{a_l},\ldots,R_{b_l}$
is the same sequence as $R_1,\ldots,R_k$ such that

- $|\dom(R_i)| = |\dom(R_j)|$ for all $m \in [1,k]$ and
all $a_m \le i \le j \le b_m$,

- $|\dom(R_{a_i})| < |\dom(R_{a_{i+1}})|$ for all $i \in [1,k-1]$.

\end{Lemma}

{\em Proof.}
By Lemma \ref{lemma:domain-stable-increase}.
$\Box$

\begin{Lemma}\label{lemma:bijection}
If $R,S$ is a domain-stable sequence
and $|\dom(S)|=|\dom(R)|$
and $R$ is injective,
then $R:\dom(R) \to \dom(S)$ is a bijection.
\end{Lemma}

{\em Proof.}
We write $R|_{X \times Y}$ for $\{ (x,y) \in X \times Y \ |\ xRy \}$.

Let $R'$ be $R|_{\dom(R) \times \dom(S)}$.

Claim 1. $\dom(R')=\dom(R)$.
We can show it as follows.
Fix $x \in \dom(R)$.
Since $R,S$ is domain-stable, we have $\dom(RS) = \dom(R)$.
Hence $x \in \dom(RS)$.
Hence there are $y$ and $z$ such that $x R y S z$.
Hence $y \in \dom(S)$.
Hence $x \in \dom(R')$.

Claim 2. $\codom(R')=\dom(S)$.
We can show it as follows.
By definition of $R'$, we have $\codom(R') \subseteq \dom(S)$.
Since $R'$ is injective, we have $|\dom(R')| \le |\codom(R')|$.
From Claim 1, we have $\dom(R')=\dom(R)$.
Hence $|\dom(R)| \le |\codom(R')| \le |\dom(S)|$.
Since $|\dom(R)| = |\dom(S)|$, we have $\codom(R')=\dom(S)$.

$R'^{-1}$ is a function $\dom(S) \to \dom(R)$.
We can show it as follows.
Fix $x \in \dom(S)$ to show $x R'^{-1} y$ for some $y$.
By Claim 2. $x \in \codom(R')$. Hence there is $y$ such that $y R' x$.
Hence $x R'^{-1} y$.
Assume $x R'^{-1} y$ and $x R'^{-1} z$ to show $y=z$.
Then $y R' x$ and $z R' x$.
Since $R'$ injective, we have $y=z$.

The function $R'^{-1}$ is surjective.
We can show it as follows.
Fix $x \in \dom(R)$.
By Claim 1, we have $x \in \dom(R')$.
Hence there is $y$ such that $x R' y$.
Hence $y R'^{-1} x$.

Since $|\dom(S)|=|\dom(R)|$, the function $R'^{-1}$ is bijective.
Hence $R'$ is a bijection.
Hence $R:\dom(R) \to \dom(S)$ is a bijection.
$\Box$

\subsection{Synthesis of ranking functions}

\begin{Lemma}\label{lemma:maxlen}
If a SC call graph satisfies SCT condition,
then
for a domain-stable injective sequence $R_1,\ldots, R_n$ of SC compositions
and a value flow $\Vec v_1,\ldots,\Vec v_n$ along $R_1,\ldots, R_n$
such that $|\dom(R_i)|=k$ for all $i \in [1,n]$ and
$\Val(R_i,\Vec v_i)$ is the same for all $i \in [1,n]$,
we have
$n \le |V|{K \choose k}$
where $V$ is the set of vertexes of the SC call graph and
$K$ is the maximum arity in all function names.
\end{Lemma}

{\em Proof.}
By Lemma \ref{lemma:bijection},
$R_i:\dom(R_i) \to \dom(R_{i+1})$ is a bijection for all $i \in [1,n-1]$.

Let $R_i=(a_i,(S_i,P_i),-,a_{i+1})$ for all $i \in [1,n]$.

Since $\Val(R_i,\Vec v_i)$ is the same for all $i \in [1,n]$,
$P_i = \emptyset$ for all $i \in [1,n-1]$.

For any $i <j \in [1,n]$, we have $\dom(R_i) \ne \dom(R_j)$ or $a_i \ne a_j$.
We can show it as follows.
Assume $\dom(R_i) = \dom(R_j)$ and $a_i = a_j$
for some $i < j \in [1,n]$ to show contradiction.
We consider an infinite composable sequence $(R_i,\ldots,R_{j-1})^\omega$.
Let $R_l$ be the composition of SC paths $R'_{l,1},\ldots,R'_{l,q_l}$
for all $l \in [i,j]$.
By applying SCT condition to
the infinite SC path 
$(R_{i,1},\dots,R_{i,q_i},\ldots,R_{j-1,1},\ldots,R_{j-1,q_{j-1}})^\omega$,
there is a progressing trace of some tail of it.
Hence
there are some $l \in [i,j-1]$ such that $R_l$ 
is progressing, which contradicts
with $P_l = \emptyset$.

Hence $(a_1,\dom(R_1)),\ldots, (a_n,\dom(R_n))$ are different.
The number of $a_i$ is $\le |V|$.
The number of $\dom(R_i)$ is $\le {K \choose k}$.
Hence $n \le |V|{K \choose k}$.
$\Box$

\begin{Def}[Weight]\label{def:weight}\rm
For a SC call graph that satisfies SCT condition,
a composable sequence $\pi$ of injective SC compositions, and 
a value flow $\sigma$ along $\pi$,
we define the weight $W(\pi, \sigma) \in ((N\cup\{\infty\}) \times N)^{K}$ 
where $K$ is the maximum arity in all function names
as follows.

Let $\pi$ be $R_1,\ldots, R_n$ and
$\sigma$ be $\Vec v_1,\ldots,\Vec v_n$.

By Proposition \ref{prop:domain-stable} (2), we define
a domain-stable sequence $R_1',\ldots,R_l'$ 
of partial SC compositions of $R_1,\ldots, R_n$ with 
head indexes $d_1,\ldots,d_l$.

We define $\Vec v'_i$ as $\Vec v_{d_i}$ for all $i \in [1,l]$.
Then by Proposition \ref{prop:domain-stable} (2), 
the value flow $\Vec v'_1,\ldots,\Vec v'_l$ is
along $R_1',\ldots,R_l'$.

By Lemma \ref{lemma:cardinal-constant},
we have
$R'_{a_1},\ldots,R'_{b_1},R'_{a_2},\ldots,R'_{b_2},\ldots, 
R'_{a_r},\ldots,R'_{b_r}$ such that

- $R'_{a_1},\ldots,R'_{b_1},R'_{a_2},\ldots,R'_{b_2},\ldots,
R'_{a_r},\ldots,R'_{b_r}$
is the same sequence as $R'_1,\ldots,R'_l$,

- $|\dom(R'_i)| = |\dom(R'_j)|$ for all $k \in [1,r]$ and
all $a_k \le i \le j \le b_k$,

- $|\dom(R'_{a_i})| < |\dom(R'_{a_{i+1}})|$ for all $i \in [1,r-1]$.

Fix $k \in [1,K]$. We define $W_k$.

Case 1. There is no $i \in [1,r]$ such that $|\dom(R'_i)| = k$.

We define $W_k$ as $(\infty,0)$.

Case 2. There is some $i \in [1,r]$ such that $|\dom(R'_i)| = k$.

Take $i \in [1,r]$ such that $|\dom(R'_{a_i})| = k$.

We define $W_k$ as $(\Val(R'_{b_i}, \Vec {v'}_{b_i}),|V|{K \choose k}-s)$
where $V$ is the set of vertexes of the SC call graph,
$K$ is the maximum arity in all function names,
and $s$ is the maximum $s$ such that
$b_i-s+1 \in [a_i,b_i]$ and
$\Val(R'_{b_i}, \Vec {v'}_{b_i}) = \Val(R'_{b_i-s+1}, \Vec {v'}_{b_i-s+1})$.
Note that by Lemma \ref{lemma:maxlen}, we have $|V|{K \choose k} \ge s$.

Note also that $s$ is the maximum length of the sequence from $R'_{b_i}$
such that each element has the same value by $\Val$.
\end{Def}

\begin{Def}\rm
We write $>_\Lex$ for the lexicographic order on $((N\cup\{\infty\}) \times N)^K$
for some $K$.
\end{Def}

\begin{Lemma}\label{lemma:stable-extension}
For a SC call graph that satisfies SCT condition,
domain-stable injective SC sequences $(\pi,\pi')$ and $(\pi,R)$ of SC compositions and 
a value flow $(\sigma,\sigma')$ along $(\pi,\pi')$ and
a value flow $(\sigma,\Vec v)$ along $(\pi,R)$,
if $\pi'=(R',-)$ and $|\dom(R)|<|\dom(R')|$ or $|\pi'|=0$,
then
\[
W((\pi,\pi'),(\sigma,\sigma'))
>_\Lex
W((\pi,R),(\sigma,\Vec v)).
\]
\end{Lemma}

{\em Proof.}

Let $\pi=(R_1,\ldots,R_n)$ and $\sigma = (\Vec v_1,\ldots,\Vec v_n)$.

Let $k=|\dom(R)|$.

Let $v = \Val(R,\Vec v)$.

Let $W((\pi,\pi'),(\sigma,\sigma')) = W = (W_1,\ldots,W_K)$ and
$W((\pi,R),(\sigma,\Vec v)) = W' = (W'_1,\ldots,W'_k)$.

We have $W_j = W'_j$ for $j \in [1,k-1]$, since $|\dom(R)|<|\dom(R')|$
or $|\pi'|=0$.

By Lemma \ref{lemma:domain-stable-increase}, 
we have $|\dom(R_{n})| \le k$.

Let $L$ be $|V|{K \choose k}$ 
where $V$ is the set of vertexes of the SC call graph
and $K$ is the maximum arity in all function names.

Case 1. $k=|\dom(R_{n})|$.

By Lemma \ref{lemma:bijection},
\[
R_n:\dom(R_n) \to \dom(R)
\]
is a bijectiion.
Hence we have $\Val(R_n,\Vec {v}_{n}) \ge v$.

Case 1.1. $\Val(R_n,\Vec {v}_{n}) = v$.

Let $W_k = (v,b)$.
Then $W'_k = (v, b-1)$.
Hence $W_k >_\Lex W'_k$.
Since $W_j = W'_j$ for $j \in [1,k-1]$,
we have $W >_\Lex W'$.

Case 1.2. $\Val(R_{n},\Vec {v}_{n}) > v$.

Let $W_k = (\Val(R_{n},\Vec {v}_{n}),b)$.
Then $W'_k = (v, L)$.
Hence $W_k >_\Lex W'_k$.
Since $W_j = W'_j$ for $j \in [1,k-1]$,
we have $W >_\Lex W'$.

Case 2. $k>|\dom(R_{n})|$.

Then $W_k = (\infty,0)$ since 
$\pi'=(R',-)$ and $|\dom(R)|<|\dom(R')|$ or $|\pi'|=0$.
We have $W'_k = (v, L)$.
Hence $W_k >_\Lex W'_k$.
Since $W_j = W'_j$ for $j \in [1,k-1]$,
we have $W >_\Lex W'$.
$\Box$

\begin{Prop}\label{prop:decrease-pi}
For a SC call graph that satisfies SCT condition, and
a composable SC sequence $(\pi,R)$ of SC compositions and 
a value flow $(\sigma,\Vec v)$ along $(\pi,R)$,
we have $W(\pi,\sigma) >_\Lex W((\pi,R),(\sigma,\Vec v))$.
\end{Prop}

{\em Proof.}
$\pi' = (\pi,R)$ and
$\sigma' = (\sigma,\Vec v)$.

Let $W(\pi,\sigma) = (W_1,\ldots,W_K)$ and
$W(\pi',\sigma') = (W'_1,\ldots,W'_K)$.

By Proposition \ref{prop:domain-stable} (2) for $\pi$, we define
a domain-stable sequence $R_1',\ldots,R_l'$ 
of partial SC compositions of $\pi$ with 
head indexes $d_1,\ldots,d_l$.

Define $\Vec v'_i$ as $\Vec v_{d_i}$ for all $i \in [1,l]$.

Then by Proposition \ref{prop:domain-stable}, 
the value flow $\Vec v'_1,\ldots,\Vec v'_l$ is
along $R'_1,\ldots,R'_l$.

Case 1. $R'_1,\ldots,R'_l,R$ is not domain-stable.

By Proposition \ref{prop:domain-stable} (1),
there is some $i$ such that $R'_1,\ldots,R'_{i-1},R'_i\ldots R'_lR$
is domain-stable and $R'_i,R'_{i+1}\ldots R'_lR$ is not domain-stable.
Let
\[
R' = R'_i\ldots R'_lR.
\]
Since $R'_i,R'_{i+1}\ldots R'_lR$ is not domain-stable and Lemma \ref{lemma:cardinal},
we have $|\dom(R')|<|\dom(R'_i)|$.

By applying Lemma \ref{lemma:stable-extension}
to $R'_1,\ldots,R'_l$ with $\Vec v'_1,\ldots,\Vec v'_l$ and
$R'_1,\ldots,R'_{i-1},R'$ with 
$\Vec v'_1,\ldots,\Vec v'_{i-1},\Vec v'_i$ with
$|\dom(R')|<|\dom(R'_i)|$,
we have
\[
W((R'_1,\ldots,R'_l),(\Vec v'_1,\ldots,\Vec v'_l)) >_\Lex
W((R'_1,\ldots,R'_{i-1},R'),
(\Vec v'_1,\ldots,\Vec v'_{i-1},\Vec v'_i)).
\]
Since
\[
W(\pi,\sigma) = W((R'_1,\ldots,R'_l),(\Vec v'_1,\ldots,\Vec v'_l)), \\
W(\pi',\sigma') = W((R'_1,\ldots,R'_{i-1},R'),
(\Vec v'_1,\ldots,\Vec v'_{i-1},\Vec v'_i))
\]
by definition, we have 
$W(\pi,\sigma) >_\Lex W(\pi',\sigma')$.

Case 2. $R'_1,\ldots,R'_l,R$ is domain-stable.

By applying Lemma \ref{lemma:stable-extension}
to $R'_1,\ldots,R'_{l}$ with $\Vec v'_1,\ldots,\Vec v'_l$ and
$R'_1,\ldots,R'_{l},R$ with $\Vec v'_1,\ldots,\Vec v'_l,\Vec v$,
we have
\[
W((R'_1,\ldots,R'_{l}), (\Vec v'_1,\ldots,\Vec v'_l)) >_\Lex
W((R'_1,\ldots,R'_{l},R), (\Vec v'_1,\ldots,\Vec v'_l,\Vec v)).
\]
Since
\[
W(\pi,\sigma) = W((R'_1,\ldots,R'_{l}), (\Vec v'_1,\ldots,\Vec v'_l)), \\
W(\pi',\sigma') = 
W((R'_1,\ldots,R'_{l},R), (\Vec v'_1,\ldots,\Vec v'_l,\Vec v))
\]
by definition, we have 
$W(\pi,\sigma) >_\Lex W(\pi',\sigma')$.
$\Box$

\begin{Th}[Intuitionistic synthesis of ranking function]\label{th:rankingfunc}
For an injective SC call graph that satisfies SCT condition,
we define
\[
\rho(f,\Vec x) = \min\{ W(\pi,\sigma) \in ((N\cup\{\infty\}) \times N)^{K} \ |\ 
\\ \qquad
\Last(\pi) = (-,-,-,f),
\sigma \hbox{\ along\ } \pi,
\Last(\sigma) \Last(\pi) \Vec x,
\pi \hbox{\ a domain-stable SC compositions},
\\ \qquad
\sigma \hbox{\ a value flow along $\pi$},
k=|\dom(\Last(\pi))|,
|\pi| \le k(|V|{K \choose k}+k!|E^+|)
\}
\]
where 
$K$ is the maximum arity in the SC call graph, and
$\Last(x)$ returns the last element of the sequence $x$, and
$V$ is the set of vertexes of the SC call graph,
and $E^+$ is the set of finite SC compositions in the call graph.

If $s \to^c s'$ for some $c \in \Callsitenames$,
then $\rho(s) >_\Lex \rho(s')$.
\end{Th}

{\em Proof.}
We define
\[
U_0(f,\Vec x) = \min\{ W(\pi,\sigma) \in ((N\cup\{\infty\}) \times N)^{K} \ |\ 
\\ \qquad
\Last(\pi) = (-,-,-,f),
\sigma \hbox{\ along\ } \pi,
\Last(\sigma) \Last(\pi) \Vec x
\},
\\
U_1(f,\Vec x) = \min\{ W(\pi,\sigma) \in ((N\cup\{\infty\}) \times N)^{K} \ |\ 
\\ \qquad
\Last(\pi) = (-,-,-,f),
\sigma \hbox{\ along\ } \pi,
\Last(\sigma) \Last(\pi) \Vec x,
\\ \qquad
\pi \hbox{\ a domain-stable sequence of SC compositions}
\}.
\]
Note that 
$U_1$ does not assume the maximum length of $\pi$,
and moreover $U_0$ does not assume the domain-stability of $\pi$.
We will show $U_0=U_1=\rho$ step by step.

Claim 1. $U_0=U_1$.

We can show it as follows.
Assume $\pi$ is not stable.
By definition of weights, we make a domain-stable $\pi'$ from $\pi$
and $\sigma'$ from $\sigma$ such that $\sigma'$ is along $\pi'$,
and $W(\pi,\rho) = W(\pi',\rho')$.
By definition of compsitions, we have $\Last(\pi')=(-,-,-,f)$
and $\Last(\sigma') \Last(\pi') \Vec x$.
Hence $U_0(f,\Vec x) \ge U_1(f,\Vec x)$.
By definition of min, we have $U_0 \le U_1$.
Hence $U_0=U_1$.

Claim 2. $U_1=\rho$.

We can show it as follows.
Fix $\pi$ and $\sigma$ such that
\[
\Last(\pi) = (-,-,-,f),
\sigma \hbox{\ along\ } \pi,
\Last(\sigma) \Last(\pi) \Vec x,
\\ \qquad
\pi \hbox{\ a domain-stable sequence of SC compositions},
\sigma \hbox{\ a value flow along $\pi$}.
\]
Let $k=|\dom(\Last(\pi))|$.

Assume $|\pi| > k(|V|{K \choose k}+k!|E^+|)$
to show that there are $\pi'$ and $\sigma'$ such that
they are obtained from $\pi$ and $\sigma$ by removing some subsequences
and $W(\pi,\sigma) = W(\pi',\sigma')$ and $|\pi|>|\pi'|$.

Let $W(\pi,\sigma)=(W_1,\ldots,W_k)$.
Since $\pi$ is domain-stable, we have $W_i=(\infty,0)$ for all $i \in [k+1,K]$.
By the condition for length, there are some $k' \in [1,k]$ and
some subsequence $R_1,\ldots,R_n$ of $\pi$
such that
for all $R \in \pi$ we have
$|\dom(R)|=k'$ iff $R=R_i$ for some $i \in [1,n]$, and
$n > |V|{K \choose k}+k!|E^+|$.
Let $\Vec v_1,\ldots,\Vec v_n$ be the subsequence of $\sigma$
corresponding to $R_1,\ldots,R_n$.
By Lemma \ref{lemma:maxlen},
there is $p \in [1,n]$ such that 
$\Val(R_i,\Vec v_i) = \Val(R_n,\Vec v_n)$ for all $i \in [p,n]$
and $\Val(R_{p-1},\Vec v_{p-1}) \ne \Val(R_p,\Vec v_p)$
and $n-p+1 \le |V|{K \choose k}$.
Hence $p-1 > k!|E^+|$.

Since $|\dom(R_i)|$ is the same for all $i \in [1,n]$,
by Lemma \ref{lemma:bijection}, $R_i : \dom(R_i) \to \dom(R_{i+1})$
is a bijection for all $i \in [1,p-1]$.
For a SC composition $R$, we write $\Bar R$ for
$\{(x,y) \in N^2 \ |\ x (S \cup P) y \}$ where $R = (-,(S,P),-,-)$.
Since the number of shapes of $R_1,\ldots,R_n$ is $\le |E^+|$,
there is some $R' \in \pi$ such that
$|\{i \in [1,p-1] \ |\ R_i=R'\}| > k!$.
Take the smallest $i \in [1,p-1]$ such that $R'=R_i$.
Since $\Bar{R_i \ldots R_j}$ is a bijection on $\dom(R')$,
we have
$|\{ \Bar{R_i \ldots R_j} \ |\ j \in [1,p-1], i \le j, R_j=R' \}| \le k!$.
Hence we have some $q,r \in [1,p-1]$ such that $i \le q<r$ and
$\Bar{R_i \ldots R_q}=\Bar{R_i \ldots R_r}$.
Since they are bijections, we have
$\Bar{R_q \ldots R_{r-1}} = 1|_{\dom(R')}$.
From bijectivity, we have $(\Vec v_q)(l) \ge (\Vec v_r)(l)$ for all $l \in \dom(R')$.

Define $\pi'$ as the sequence obtained from $\pi$ by removing
$R_{q},\ldots,R_{r-1}$ and 
define $\Hat\sigma$ as the sequence obtained from $\sigma$
corresponding to $\pi'$.
Then $|\pi'|<|\pi|$ and $\pi'$ is domain-stable.
Moreover $\Hat\sigma$ is along $\pi'$, since $\Vec v_q \ge \Vec v_r$ on $\dom(R')$
and $\codom(R_{q-1}) \subseteq \dom(R_q)$ from the domain-stability.

Then by comparing the $i$-th element for each $i \in [1,K]$,
we have
$W(\pi',\Hat\sigma)=W(\pi,\sigma)$.
We have shown Claim 2.

By claims 1 and 2, we have $\rho=U_0$.

Assume $(f,\Vec x) \to^c (g,\Vec v)$ to show
$U_0(f,\Vec x) >_\Lex U_0(g,\Vec v)$.

Let $R$ be the edge $(f,-,c,g)$.

Define $\pi$ and $\sigma$ by
\[
U_0(f,\Vec x)=W(\pi,\sigma), \\
\Last(\pi) = (-,-,-,f), \\
\sigma \hbox{\ along\ } \pi,
\Last(\sigma) \Last(\pi) \Vec x.
\]
Then 
$(\sigma,\Vec v)$ is along $(\pi,R)$.
By Proposition \ref{prop:decrease-pi},
we have
$W(\pi,\sigma) >_\Lex W((\pi,R),(\sigma,\Vec v))$.
Hence
$U_0(f,\Vec x) =
W(\pi,\sigma) >_\Lex W((\pi,R),(\sigma,\Vec v))
\ge_\Lex U_0(g,\Vec v)$.
Hence $U_0(f,\Vec x) >_\Lex U_0(g,\Vec v)$.
$\Box$

Note that the ranking function synthesis in this theorem
is done in intuitionistic logic.

Note also that the upper bound $k(|V|{K \choose k}+k!|E^+|)$ for $|\pi|$
can be overapproximated by the SC graph as lollows.
$|E^+| \le |V|^2C2^{2K^2}$ from $S,P \subseteq (2^K)^K$
and ${K \choose k} \le 2^K$
and $k \le K$ where $C = |\Callsitenames|$.

